% Part: incompleteness
% Chapter: incompleteness-provability
% Section: second-incompleteness-thm

\documentclass[../../../include/open-logic-section]{subfiles}

\begin{document}

\olfileid{inc}{inp}{2in}

\olsection{दुसरे अपूर्णता प्रमेय}

$\Th{PA}$ स्वतःची सुसंगतता सिद्ध करू
शकत नाही, हा दावा कसा व्यक्त करायचा?
$\Th{PA}$ विसंगत आहे असे म्हणणे
म्हणजे $\Th{PA} \Proves \eq[0][1]$
असे म्हणणे. सुसंगतता-विधानासाठी
असत्यता-स्थिराचा संकेतांक घेतला तरी
सुसंगततेचा आशय तोच राहतो; या
सिद्धतेत $\OCon[\Th{PA}]$ हे
!!{sentence} $\lnot
\OProv[\Th{PA}](\gn{\lfalse})$
मानू. मग पुढील प्रमेय अपेक्षित
निकाल देते:

\begin{thm}
\ollabel{thm:second-incompleteness}
$\Th{PA}$ सुसंगत असेल, तर
$\Th{PA}$ हे
$\OCon[\Th{PA}]$ !!{derive}
करत नाही.
\end{thm}

हे प्रमेय $\OCon[\Th{PA}]$
च्या विशिष्ट निरूपणावर
(म्हणजे $\OProv[\Th{PA}](y)$
च्या विशिष्ट निरूपणावर)
अवलंबून आहे, हे लक्षात
घेणे महत्त्वाचे आहे.
आपण फक्त एवढेच वापरणार
आहोत की $\OProv[\Th{PA}](y)$
चे निरूपण तीन
!!{derivability} अटी पूर्ण
करते. म्हणून असे गुणधर्म
असलेले !!a{derivability}
विधेय असलेल्या कोणत्याही
उपपत्तीला हे प्रमेय लागू
होते.

गोडेलच्या युक्तिवादाची
रूपरेषा वाचणे उपयुक्त आहे;
प्रमेय त्यातून नेमके
समारोपाप्रमाणे निघते.
युक्तिवाद असा आहे.
\olref[1in]{thm:first-incompleteness}
च्या सिद्धतेत आपण तयार
केलेले गोडेल-वाक्य
$!G_\Th{PA}$ असू द्या.
“$\Th{PA}$ सुसंगत
असेल, तर $\Th{PA}$
हे $!G_\Th{PA}$ ला
!!{derive} करत नाही”
हे आपण दाखवले आहे.
हे $\Th{PA}$
\emph{मध्येच} आकारिक
केल्यास पुढील विधानाची
सिद्धता मिळते:
\[
\OCon[\Th{PA}] \lif \lnot \OProv[\Th{PA}](\gn{!G_\Th{PA}}).
\]
आता $\Th{PA}$ हे
$\OCon[\Th{PA}]$
!!{derive}s करते असे
मानू. मग ते $\lnot
\OProv[\Th{PA}](\gn{!G_\Th{PA}})$
सुद्धा !!{derive}s
करते. पण $!G_\Th{PA}$
हे गोडेल-वाक्य असल्यामुळे
हे $!G_\Th{PA}$ शी
सममूल्य आहे. म्हणून
$\Th{PA}$ हे
$!G_\Th{PA}$
!!{derive}s करते.

पण $\Th{PA}$ सुसंगत
असेल, तर ते
$!G_\Th{PA}$ ला
!!{derive} करत
नाही, हे आपल्याला
माहीत आहे!{}
म्हणून $\Th{PA}$
सुसंगत असेल, तर
$\OCon[\Th{PA}]$
सुद्धा !!{derive}
करू शकत नाही.

युक्तिवाद अधिक नेमका
करण्यासाठी $!G_\Th{PA}$
हे~$\Th{PA}$ चे
गोडेल-वाक्य मानू.
!!{derivability} अटी
(P1)--(P3) वापरून
$\Th{PA}$ हे
$\OCon[\Th{PA}] \lif
!G_\Th{PA}$ !!{derive}s
करते हे दाखवू.
त्यावरून $\Th{PA}$
हे $\OCon[\Th{PA}]$
!!{derive} करत नाही
हे मिळेल. पुढे
$\Th{PA}$ मधील
सिद्धतेची रूपरेषा आहे.
(सोयीसाठी
$\Th{PA}$ हे अधोलिखित
वगळू.)
\begin{align}
& !G \liff \lnot \OProv(\gn{!G}) \ollabel{G2-1}\\
& \qquad\text{$!G$ हे गोडेल-वाक्य आहे}\notag \\
& !G \lif \lnot \OProv(\gn{!G}) \ollabel{G2-2}\\
  & \qquad\text{\olref{G2-1} वरून} \notag\\
& !G \lif
  (\OProv(\gn{!G}) \lif \lfalse) \ollabel{G2-3}\\
  & \qquad\text{\olref{G2-2} व तर्कशास्त्रावरून}\notag\\
& \OProv(\gn{
    !G \lif
    (\OProv(\gn{!G}) \lif \lfalse)
  }) \ollabel{G2-4}\\
  & \qquad\text{\olref{G2-3} आणि अट P1 वरून} \notag\\
& \OProv(\gn{!G}) \lif
  \OProv(\gn{
    (\OProv(\gn{!G}) \lif \lfalse)
    }) \ollabel{G2-5}\\
  & \qquad\text{\olref{G2-4} आणि अट P2 वरून} \notag\\
& \OProv(\gn{!G}) \lif (\OProv(\gn{\OProv(\gn{!G})}) \lif \OProv(\gn{\lfalse})) \ollabel{G2-6}\\
  & \qquad\text{\olref{G2-5}, अट P2 आणि तर्कशास्त्रावरून} \notag\\
& \OProv(\gn{!G}) \lif
  \OProv(\gn{\OProv(\gn{!G})}) \ollabel{G2-7}\\
   & \qquad\text{अट P3 वरून} \notag\\
& \OProv(\gn{!G}) \lif \OProv(\gn{\lfalse}) \ollabel{G2-8}\\
  & \qquad \text{\olref{G2-6}, \olref{G2-7} व तर्कशास्त्रावरून}\notag\\
& \OCon \lif \lnot \OProv(\gn{!G}) \ollabel{G2-9}\\
  & \qquad\text{\olref{G2-8} च्या प्रतिवर्ती विधानावरून आणि $\OCon \ident \lnot \OProv(\gn{\lfalse})$ वरून}\notag \\
& \OCon \lif !G \notag\\
  & \qquad\text{\olref{G2-1}, \olref{G2-9} व तर्कशास्त्रावरून}\notag
\end{align}
वरील तर्कशास्त्राचा
वापर केवळ विधानीय
तर्कशास्त्रातील प्राथमिक
तथ्यांपुरता आहे.
उदाहरणार्थ,
\olref{G2-3} मध्ये
$\Proves \lnot!A \liff (!A\lif
\lfalse)$ वापरले आहे,
आणि \olref{G2-8} मध्ये
$!A \lif (!B \lif !C), !A \lif !B
\Proves !A \lif !C$ वापरले
आहे. \olref{G2-5} आणि
\olref{G2-6} मध्ये अट~P2
वापरताना तिची पुढील
उदाहरणे वापरली आहेत:
$\OProv(\gn{!A \lif !B}) \lif
(\OProv(\gn{!A}) \lif \OProv(\gn{!B}))$.
पहिल्यात $!A \ident
!G$ आणि $!B \ident
\OProv(\gn{!G}) \lif \lfalse$;
दुसऱ्यात $!A
\ident \OProv(\gn{!G})$
आणि $!B \ident \lfalse$.

दुसऱ्या अपूर्णता प्रमेयाचे
अधिक अमूर्त रूप असे:

\begin{thm}
\ollabel{thm:second-incompleteness-gen}
$\Th{T}$ ही~$\Th{Q}$ चा
विस्तार असलेली कोणतीही
सुसंगत, !!{axiomatized}
उपपत्ती असू द्या.
$\OProv[\Th{T}](y)$
हे सूत्र $\Th{T}$
साठीच्या
!!{derivability} अटी
P1--P3 पूर्ण करत असेल,
तर $\Th{T}$ हे
$\OCon[\Th{T}]$
!!{derive} करत नाही.
\end{thm}

\begin{prob}
$\Th{PA}$ हे
$!G_{\Th{PA}} \lif \OCon[\Th{PA}]$
!!{derive}s करते हे
दाखवा.
\end{prob}

\begin{digress}
या कथेतून असा बोध होतो:
गणितासाठीची कोणतीही
“स्वीकारार्ह” सुसंगत
उपपत्ती स्वतःचे
सुसंगतता-विधान
!!{derive} करू शकत
नाही. $\Th{T}$ ही
गणिताची उपपत्ती असून
तिच्यात $\Th{Q}$ आणि
हिल्बर्टचे “सांतवादी”
युक्तिवाद (त्यांचा नेमका
अर्थ काहीही असो)
समाविष्ट आहेत असे
मानू. मग संपूर्ण
$\Th{T}$ हे
$\Th{T}$ चे
सुसंगतता-विधान
!!{derive} करू
शकत नाही. त्यामुळे
तिचा सांतवादी भागही
$\Th{T}$ चे
सुसंगतता-विधान
!!{derive} करू
शकत नाही, हे तर
अधिकच स्पष्ट आहे.
त्या अर्थाने
“सर्व गणितासाठी”
सुसंगततेची सांतवादी
सिद्धता असू शकत नाही.

“सांतवादी” या संज्ञेचा
अर्थ लावण्यात थोडी
मोकळीक आहे. गोडेलने
1931 च्या शोधपत्रात
हिल्बर्टला आकारिक
करायचे असलेल्या
गणिताबाहेरही आपण
“सांतवादी” म्हणू
असे काही असू शकेल,
ही शक्यता मान्य केली.
पण ही गोडेलची
उदार भूमिका होती:
आज सांतवादी म्हणता
येईल पण, उदाहरणार्थ,
निवडीच्या स्वयंसिद्धकासह
झर्मेलो--फ्रेंकेल संच
उपपत्ती $\Th{ZFC}$
मध्ये आकारिक करता
येणार नाही असे
काही सापडेल, हे
कल्पना करणे कठीण आहे.
\end{digress}

\end{document}
